\documentclass[12pt, a4paper]{article}

% Packages
\usepackage[margin=1in]{geometry}
\usepackage{setspace}
\usepackage{parskip}
\usepackage{titlesec}
\usepackage{fancyhdr}
\usepackage{listings}
\usepackage{xcolor}
\usepackage{hyperref}
\usepackage{microtype}
\usepackage{apacite}

% Font
\usepackage{times}

% Double spacing (APA)
\doublespacing

% Header / footer
\pagestyle{fancy}
\fancyhf{}
\rhead{Function Interface and Streams}
\lhead{Dau, A.}
\rfoot{\thepage}
\renewcommand{\headrulewidth}{0.4pt}

% Section formatting
\titleformat{\section}{\normalfont\bfseries\centering}{}{0em}{}
\titleformat{\subsection}{\normalfont\bfseries}{}{0em}{}

% Code listing style
\definecolor{codegray}{rgb}{0.95,0.95,0.95}
\definecolor{commentgreen}{rgb}{0.1,0.5,0.1}
\definecolor{keywordblue}{rgb}{0.1,0.1,0.7}
\definecolor{stringred}{rgb}{0.6,0.1,0.1}

\lstdefinestyle{java}{
	backgroundcolor=\color{codegray},
	commentstyle=\color{commentgreen}\itshape,
	keywordstyle=\color{keywordblue}\bfseries,
	stringstyle=\color{stringred},
	basicstyle=\ttfamily\footnotesize\singlespacing,
	breakatwhitespace=false,
	breaklines=true,
	keepspaces=true,
	numbers=left,
	numbersep=8pt,
	numberstyle=\tiny\color{gray},
	showspaces=false,
	showstringspaces=false,
	tabsize=4,
	frame=single,
	framesep=4pt,
	language=Java,
	captionpos=b
}

\lstdefinestyle{output}{
	backgroundcolor=\color{codegray},
	basicstyle=\ttfamily\footnotesize\singlespacing,
	breaklines=true,
	frame=single,
	framesep=4pt,
	numbers=none
}

\lstset{style=java}

\hypersetup{colorlinks=true, linkcolor=black, urlcolor=black, citecolor=black}

% ---------------------------------------------------------------
\begin{document}
	
	% Title page
	\begin{titlepage}
		\centering
		\vspace*{2cm}
		{\large University of the People \par}
		\vspace{0.5cm}
		{\large CS-1102: Programming 1 \par}
		\vspace{3cm}
		{\LARGE\bfseries The \texttt{Function} Interface and Streams in Java \par}
		\vspace{3cm}
		{\large Andreas Dau \par}
		{\large Student ID: C644260 \par}
		\vspace{1cm}
		{\large Programming Assignment, Unit 8 \par}
		\vspace{1cm}
		{\large June 4, 2026 \par}
		\vfill
	\end{titlepage}
	
	% ---------------------------------------------------------------
	\section{Background}
	
	Java 8 introduced functional interfaces and the Stream API to allow developers
	to express data transformations more directly \cite{liang2020}. The
	\texttt{Function<T,\,R>} interface represents any operation that takes one input
	of type \texttt{T} and produces one output of type \texttt{R}. Its single
	abstract method, \texttt{apply(T t)}, can be implemented with a lambda expression,
	making it possible to store, pass, and combine functions as values
	\cite{urma2018}. The related \texttt{Predicate<T>} interface represents a
	boolean-valued function and supports composition through \texttt{and()},
	\texttt{or()}, and \texttt{negate()}. Streams build on these interfaces: a
	stream pipeline chains intermediate operations such as \texttt{filter()} and
	\texttt{map()}, which accept functional interfaces as arguments, and defers
	execution until a terminal operation such as \texttt{toList()} or
	\texttt{average()} is called \cite{goetz2006}. This lazy evaluation model keeps
	memory consumption bounded even for large datasets.
	
	% ---------------------------------------------------------------
	\section{Program Description}
	
	The program is a command-line tool that processes a dataset of company
	employees. Each employee has four attributes: name, age, department, and salary.
	The tool accepts one or more \texttt{--filter} expressions at runtime, builds a
	composite \texttt{Predicate<Employee>} from them, and passes it to a stream
	pipeline that produces a filtered list of names and departments alongside the
	average salary of the matching employees.
	
	\subsection{The \texttt{Function} Interface in Use}
	
	A \texttt{Function<Employee, String>} is defined to concatenate each employee's
	name and department into a single string. This function is passed directly to
	\texttt{Stream.map()}, which applies it to every element in the filtered stream:
	
	\begin{lstlisting}
		Function<Employee, String> nameAndDept =
		emp -> emp.getName() + ", " + emp.getDepartment();
		
		List<String> result = employees.stream()
		.filter(combinedFilter)
		.map(nameAndDept)
		.toList();
	\end{lstlisting}
	
	\subsection{Filter Expressions and \texttt{Predicate} Composition}
	
	The \texttt{FilterParser} class parses a filter string such as \texttt{age>30}
	or \texttt{department=Engineering} into a \texttt{Predicate<Employee>}. The
	supported operators are \texttt{=}, \texttt{\textasciitilde} (contains),
	\texttt{>}, \texttt{<}, \texttt{>=}, and \texttt{<=}. When multiple
	\texttt{--filter} flags are given, the predicates are combined with
	\texttt{Predicate.and()}, so all conditions must match:
	
	\begin{lstlisting}
		Predicate<Employee> combined = emp -> true;
		for (int i = 0; i < args.length - 1; i++) {
			if ("--filter".equals(args[i])) {
				combined = combined.and(FilterParser.parse(args[i + 1]));
			}
		}
	\end{lstlisting}
	
	\subsection{Average Salary}
	
	The average salary is computed using \texttt{mapToDouble()} to produce a
	primitive \texttt{DoubleStream}, followed by the built-in \texttt{average()}
	reduction. The result is wrapped in \texttt{OptionalDouble} to handle an
	empty dataset without throwing an exception:
	
	\begin{lstlisting}
		OptionalDouble avgSalary = employees.stream()
		.filter(combinedFilter)
		.mapToDouble(Employee::getSalary)
		.average();
		
		avgSalary.ifPresentOrElse(
		avg -> System.out.printf("$%.2f%n", avg),
		()  -> System.out.println("(no data)"));
	\end{lstlisting}
	
	\subsection{CSV Support and \texttt{EmployeeFactory}}
	
	The \texttt{-i}/\texttt{--infile} flag loads a CSV file instead of the
	built-in dataset. The companion tool \texttt{EmployeeFactory} generates CSV
	files of arbitrary size for testing:
	
	\begin{lstlisting}[style=output]
		java EmployeeFactory 1000000 big_dataset.csv
		java EmployeeProcessor -i big_dataset.csv --filter age>30 --filter salary>=70000
	\end{lstlisting}
	
	The stream pipeline requires no modification for larger datasets because
	elements are processed one at a time through lazy evaluation.
	
	\subsection{Design Decision: Class Instead of Record}
	
	The \texttt{Employee} data model is implemented as a static inner class rather
	than a Java record. While records, stable since Java 16, would generate the
	constructor and accessors automatically, the explicit class form was chosen to
	keep the structure transparent and consistent with the course material.
	
	% ---------------------------------------------------------------
	\section{Source Code}
	
	\subsection{\texttt{Employee.java}}
	
	\lstinputlisting{../src/Employee.java}
	
	\subsection{\texttt{FilterParser.java}}
	
	\lstinputlisting{../src/FilterParser.java}
	
	\subsection{\texttt{CSVReader.java}}
	
	\lstinputlisting{../src/CSVReader.java}
	
	\subsection{\texttt{EmployeeFactory.java}}
	
	\lstinputlisting{../src/EmployeeFactory.java}
	
	\subsection{\texttt{EmployeeProcessor.java}}
	
	\lstinputlisting{../src/EmployeeProcessor.java}
	
	% ---------------------------------------------------------------
	\section{Sample Output}
	
	\subsection{No filters}
	
	\begin{lstlisting}[style=output]
		> java EmployeeProcessor
		=== Active filters ===
		(none - showing all employees)
		
		=== Name, Department ===
		Alice, Engineering
		Bob, Marketing
		Carol, Engineering
		David, HR
		Eve, Marketing
		Frank, Finance
		
		=== Average Salary ===
		$73000.00
	\end{lstlisting}
	
	\subsection{Filter: age above 30 and salary at least \$70,000}
	
	\begin{lstlisting}[style=output]
		> java EmployeeProcessor --filter age>30 --filter salary>=70000
		=== Active filters ===
		age>30
		salary>=70000
		
		=== Name, Department ===
		Alice, Engineering
		Carol, Engineering
		Frank, Finance
		
		=== Average Salary ===
		$91000.00
	\end{lstlisting}
	
	\subsection{Filter: department contains ``ing''}
	
	\begin{lstlisting}[style=output]
		> java EmployeeProcessor --filter department~ing
		=== Active filters ===
		department~ing
		
		=== Name, Department ===
		Alice, Engineering
		Carol, Engineering
		Frank, Finance
		
		=== Average Salary ===
		$91000.00
	\end{lstlisting}
	
	\subsection{Help output}
	
	\begin{lstlisting}[style=output]
		> java EmployeeProcessor --help
		Usage:
		java EmployeeProcessor [-i file] [--filter expression] ...
		
		Options:
		-h, --help            Show this help message
		-?, --help            Same as above
		-i, --infile file     Read employees from a CSV file instead of
		the built-in dataset
		--filter expr         Filter employees by a field expression (repeatable)
		
		Filter expressions:
		field operator value
		
		Fields:     name, age, department, salary
		Operators:  =   exact match (all fields)
		~   contains, case-insensitive (name, department)
		>   greater than (age, salary)
		<   less than (age, salary)
		>=  greater than or equal (age, salary)
		<=  less than or equal (age, salary)
		
		Multiple --filter flags are combined with AND logic.
	\end{lstlisting}
	
	\subsection{Error handling}
	
	\begin{lstlisting}[style=output]
		> java EmployeeProcessor --filter age>abc
		Error in filter 'age>abc': 'abc' is not a valid number for field 'age'
		Run with -h or --help for usage information.
		
		> java EmployeeProcessor --badflag
		Error: Unknown option '--badFlag'
		Run with -h or --help for usage information.
	\end{lstlisting}
	
	% ---------------------------------------------------------------
	\newpage
	\singlespacing
	\bibliographystyle{apacite}
	\bibliography{references}
	
\end{document}
